% Pista alg-26s-q2 · Àlgebra
% Procedència: PAU Catalunya 2026 · Convocatòria de setembre · Sèrie 2 · Exercici 2
\documentclass[11pt,a4paper]{article}
\usepackage{pista}
\pistainfo{Catalunya 2026}{Convocatòria de setembre}{Sèrie 2}

\begin{document}

\section*{\textcolor{blauPista}{Exercici 2}}

\noindent Considereu el sistema d'equacions següent:
\[
\left.\begin{array}{r}
2x + ky + z = 3 \\
-x + 2y + 2z = 1 \\
x + y + z = 2
\end{array}\right\},
\]
on $k$ és un paràmetre real.

\subsection*{2.1. \normalfont Discutiu el sistema en funció del paràmetre $k$. \hspace*{\fill}\textnormal{[1~punt]}}

\pista{Escriu la matriu $A$ i la matriu ampliada $A^{'}$, i calcula el determinant de la matriu $A$, aplicant la regla de Sarrus.

\medskip
A continuació, estudia el sistema d'equacions. D'una banda, has d'estudiar a part cada valor del paràmetre $k$ que anul·la el determinant de la matriu $A$. D'altra banda, la resta de valors que no anul·len el determinant fan que el sistema sigui compatible determinat, és a dir, que té solució única.

\medskip
Per a cada valor del paràmetre $k$ que anul·la el determinant, busca un menor d'ordre 2 de la matriu $A$ que no valgui zero. A continuació, per trobar el rang de la matriu $A^{'}$, orla'l amb la fila que hi falta i amb la columna dels termes independents.}

\subsection*{2.2. \normalfont Resoleu el sistema per a tots els valors de $k$ i comproveu que, per a $k \neq 1$, la solució no depèn de $k$. \hspace*{\fill}\textnormal{[1~punt]}}

\pista{Esglaona la matriu ampliada $A^{'}$ de l'apartat anterior. En algun moment obtindràs una equació semblant a aquesta:
\[
(1 - k) \cdot y = 0
\]
Si $k \neq 1$, podràs trobar un valor concret per a $x$, $y$ i $z$, i veuràs per què la solució no depèn de $k$.

\medskip
En canvi, per a $k = 1$, tindràs una equació que es compleix sempre. Llavors, escriu simplement \mbox{$y = t$} (és un paràmetre lliure), i expressa les altres incògnites en funció de $t$.}

\subsection*{2.3. \normalfont Interpreteu geomètricament el resultat dels apartats anteriors. \hspace*{\fill}\textnormal{[0,5~punts]}}

\pista{Cadascuna de les equacions s'interpreta geomètricament com un pla. La 2a equació i la 3a equació no depenen de $k$: són dos plans fixos, que es tallen en una recta. En canvi, el pla de la 1a equació depèn del paràmetre $k$.

\medskip
Per a $k \neq 1$, on es tallen els tres plans? I per a $k = 1$: creus que es tallen en un punt, en una recta, o bé que no tenen cap punt en comú? Si es tallen en una recta, quina ha de ser?

\medskip
Finalment, pensa què vol dir que la solució per a $k \neq 1$ no depengui de $k$: per quin punt passen tots els plans de la 1a equació?}

\end{document}
